% !PRINTED 562
\weberpage{562}
%<<CONT
\weq{9}{\begin{aligned}
\eta &= r + r^{-5} + r^{-1} + r^5 = \eta^{(1)}\\
\eta_1 &= r^2 + r^3 + r^{-2} + r^{-3} = \eta^{(2)}\\
\eta_2 &= r^4 + r^6 + r^{-4} + r^{-6} = \eta^{(4)}
\end{aligned}}

If one applies formula (4), one obtains, e.g.,

\[
\begin{aligned}
\eta\eta &= \eta^{(2)} + \eta^{(-4)} + \eta^{(0)} + \eta^{(6)} = 4 + \eta^{(2)} + 2\eta^{(4)}\\
&= -4\eta - 3\eta_1 - 2\eta_2,
\end{aligned}
\]

and thus the formulae result

\[
\begin{aligned}
\eta^2 &= -4\eta - 3\eta_1 - 2\eta_2,\\
\eta\eta_1 &= \eta + 2\eta_1 + \eta_2,\\
\eta\eta_2 &= 2\eta + \eta_1 + \eta_2,
\end{aligned}
\]

from which, by \S.\,167, (17), the equation for $\eta$

\[
\begin{vmatrix}
-4-\eta, & -3, & -2\\
1, & 2-\eta, & 1\\
2, & 1, & 1-\eta
\end{vmatrix} = 0,
\]

or

\weq{10}{\eta^3 + \eta^2 - 4\eta + 1 = 0}

follows. The discriminant of this equation is found to be equal to
$169 = 13^2$, hence positive. The equation consequently has three real
% [sic] The print attributes the sign pattern to the positive constant term alone. That
% term shows only that an odd number of roots is negative; the negative coefficient of
% eta supplies the additional exclusion of three negative roots. The conclusion is correct.
roots, and indeed, since the last term is positive, two positive and one
negative root.

If we set $r = e^{\frac{2\pi i}{13}}$, then

\[
\begin{aligned}
\eta &= 2\cos\frac{2\pi}{13} + 2\cos\frac{10\pi}{13} = 2\left(\cos\frac{2\pi}{13} - \cos\frac{3\pi}{13}\right),\\
\eta_1 &= 2\cos\frac{4\pi}{13} + 2\cos\frac{6\pi}{13} = 2\left(\cos\frac{4\pi}{13} + \cos\frac{6\pi}{13}\right),\\
\eta_2 &= 2\cos\frac{8\pi}{13} + 2\cos\frac{12\pi}{13} = -2\left(\cos\frac{\pi}{13} + \cos\frac{5\pi}{13}\right),
\end{aligned}
\]

or also

\[
\begin{aligned}
\eta &= 4\cos\frac{4\pi}{13}\cos\frac{6\pi}{13}, \qquad \eta_1 = 4\cos\frac{\pi}{13}\cos\frac{5\pi}{13},\\
\eta_2 &= -4\cos\frac{2\pi}{13}\cos\frac{3\pi}{13}.
\end{aligned}
\]

Thus $\eta_2$ is the negative root, $\eta_1$ the greater, and $\eta$ the smaller of the
two positive roots.

The equation of degree $4$, whose roots are $r$, $r^{-1}$, $r^5$, $r^{-5}$,
can now easily be formed, if one
%CONT>>
